\section*{第 \ref{ch:b1:complex} 章——复数}

\begin{solution}{exo:b1:complex:1}
$1 + \iu = \sqrt 2\, \eu^{\iu\pi/4}$；
$\;\sqrt 3 - \iu = 2\, \eu^{-\iu\pi/6}$；
$\;-5 = 5\, \eu^{\iu\pi}$；
\[
\frac{1 + \iu}{\sqrt 3 - \iu}
= \frac{\sqrt 2}{2}\, \eu^{\iu(\pi/4 + \pi/6)}
= \frac{\sqrt 2}{2}\, \eu^{5\iu\pi/12}.
\]
代数地计算同一商（分子分母同乘分母的\omterm{def:b1:complex:field}{共轭复数}）：
\[
\frac{(1 + \iu)(\sqrt 3 + \iu)}{4}
= \frac{(\sqrt 3 - 1) + \iu(\sqrt 3 + 1)}{4}.
\]
将其与 $\frac{\sqrt 2}{2}(\cos\frac{5\pi}{12} +
\iu\sin\frac{5\pi}{12})$ 的三角形式比较，得：
\[
\cos\frac{5\pi}{12} = \frac{\sqrt 6 - \sqrt 2}{4},
\qquad
\sin\frac{5\pi}{12} = \frac{\sqrt 6 + \sqrt 2}{4}.
\]
\end{solution}

\begin{solution}{exo:b1:complex:2}
$(1+\iu)^2 = 2\iu$，故 $(1+\iu)^{20} = (2\iu)^{10} = 2^{10}\,\iu^{10}
= 1024 \times (-1) = -1024$。

由指数形式 $(1+\iu)^n = 2^{n/2}\, \eu^{\iu n\pi/4}$，
它为实数当且仅当 $\sin\frac{n\pi}{4} = 0$，即 $4 \mid n$。故 $(1+\iu)^n
\in \R$ 恰对 $4$ 的倍数成立（此时取值 $(-4)^{n/4}$）。
\end{solution}

\begin{solution}{exo:b1:complex:3}
$\abs{z - 2} = \abs{z + \iu}$：到点 $2$ 与
$-\iu$ 等距的点——即连接 $(2, 0)$
与 $(0, -1)$ 的线段的垂直平分线。

$\abs{z - 1} = 2$：圆心为 $1$、半径为 $2$ 的圆。

$\frac{z-1}{z+1} \in \iu\R$：由
\cref{prop:b1:complex:geometry}~(3)，点 $M(z)$、$A(1)$
、$B(-1)$ 满足：直线 $MA$ 与 $MB$ 垂直（或 $z = 1$，此时
商为 $0 \in \iu\R$）。所求\omterm{def:b1:logic:sets}{集合}是
以 $[-1, 1]$ 为直径的圆（即单位圆），去掉点 $-1$，该点处的商
无定义。\emph{用计算验证：}$z = \eu^{\iu\theta}$ 得
$\frac{z-1}{z+1} = \frac{\eu^{\iu\theta/2}(\eu^{\iu\theta/2} -
\eu^{-\iu\theta/2})}{\eu^{\iu\theta/2}(\eu^{\iu\theta/2} +
\eu^{-\iu\theta/2})} = \iu\tan\frac\theta2 \in \iu\R$。
\end{solution}

\begin{solution}{exo:b1:complex:4}
线性化：
\[
\sin^4\theta
= \Bigl(\frac{\eu^{\iu\theta} - \eu^{-\iu\theta}}{2\iu}\Bigr)^{\!4}
= \frac{\eu^{4\iu\theta} - 4\eu^{2\iu\theta} + 6 - 4\eu^{-2\iu\theta}
+ \eu^{-4\iu\theta}}{16}
= \frac{\cos 4\theta - 4\cos 2\theta + 3}{8}.
\]
展开：由棣莫弗公式，$\cos 4\theta = \Re\bigl((c +
\iu s)^4\bigr) = c^4 - 6c^2 s^2 + s^4$，其中 $c = \cos\theta$，$s =
\sin\theta$；代入 $s^2 = 1 - c^2$，得：
\[
\cos 4\theta = c^4 - 6c^2(1 - c^2) + (1 - c^2)^2
= 8c^4 - 8c^2 + 1 .
\]
\end{solution}

\begin{solution}{exo:b1:complex:5}
$8\iu = 8\,\eu^{\iu\pi/2}$，故立方根为 $2\,\eu^{\iu(\pi/6 +
2k\pi/3)}$，$k = 0, 1, 2$：
\[
z_0 = 2\eu^{\iu\pi/6} = \sqrt 3 + \iu,\quad
z_1 = 2\eu^{5\iu\pi/6} = -\sqrt 3 + \iu,\quad
z_2 = 2\eu^{3\iu\pi/2} = -2\iu :
\]
这三点构成半径为 $2$ 的圆上的正三角形。

$-4 = 4\eu^{\iu\pi}$，故 $z^4 = -4$ 的解为 $\sqrt 2\,
\eu^{\iu(\pi/4 + k\pi/2)}$：$\;1 + \iu$、$-1 + \iu$、$-1 - \iu$、
$1 - \iu$。将共轭根两两配对：
\[
X^4 + 4 = \bigl(X^2 - 2X + 2\bigr)\bigl(X^2 + 2X + 2\bigr),
\]
由于 $(X - (1+\iu))(X - (1-\iu)) = X^2 - 2X + 2$，另一组
类似。（展开即可验证。）
\end{solution}

\begin{solution}{exo:b1:complex:6}
$C_n + \iu S_n = \sum_{k=0}^{n} \eu^{\iu k\theta} =
\frac{\eu^{\iu(n+1)\theta} - 1}{\eu^{\iu\theta} - 1}$（等比求和，
公比 $\eu^{\iu\theta} \neq 1$）。半角因式分解
（见\cref{met:b1:complex:trig}~(4)），作用于分子与分母：
\[
\frac{2\iu\sin\frac{(n+1)\theta}{2}\, \eu^{\iu(n+1)\theta/2}}
{2\iu\sin\frac{\theta}{2}\, \eu^{\iu\theta/2}}
= \frac{\sin\frac{(n+1)\theta}{2}}{\sin\frac{\theta}{2}}\,
\eu^{\iu n\theta/2}.
\]
分别取实部与虚部：
\[
C_n = \frac{\sin\frac{(n+1)\theta}{2}}{\sin\frac{\theta}{2}}
\cos\frac{n\theta}{2},
\qquad
S_n = \frac{\sin\frac{(n+1)\theta}{2}}{\sin\frac{\theta}{2}}
\sin\frac{n\theta}{2}.
\]
\end{solution}

\begin{solution}{exo:b1:complex:7}
判别式：$\Delta = (3 + 4\iu)^2 - 4(-1 + 5\iu) = 9 + 24\iu - 16 +
4 - 20\iu = -3 + 4\iu$。$-3 + 4\iu$ 的平方根：解 $x^2 - y^2
= -3$，$2xy = 4$，$x^2 + y^2 = 5$；于是 $x^2 = 1$，$y = 2/x$，二者同
号：$\delta = \pm(1 + 2\iu)$。由此得
\[
z = \frac{(3 + 4\iu) \pm (1 + 2\iu)}{2}
\in \{\, 2 + 3\iu,\; 1 + \iu \,\}.
\]
\emph{检验：}和 $= 3 + 4\iu$，积 $(2+3\iu)(1+\iu) = -1 +
5\iu$，与系数相符。
\end{solution}

\begin{solution}{exo:b1:complex:8}
\begin{enumerate}
  \item 在\cref{prop:b1:complex:sumroots}中取 $n = 5$。
  \item 由 (1)，$u + v = \omega + \omega^2 + \omega^3 + \omega^4 = -1$。
        求乘积时展开，并将指数模 $5$ 约化：
        \[
        uv = (\omega + \omega^4)(\omega^2 + \omega^3)
        = \omega^3 + \omega^4 + \omega^6 + \omega^7
        = \omega^3 + \omega^4 + \omega + \omega^2 = -1 .
        \]
        故 $u$、$v$ 的和为 $-1$，积为 $-1$，它们是
        $X^2 + X - 1$ 的两个根。
  \item $u = \omega + \conj\omega = 2\cos\frac{2\pi}{5} > 0$（该角为
        锐角），且 $X^2 + X - 1$ 的正根为
        $\frac{-1 + \sqrt 5}{2}$。由此得 $\cos\frac{2\pi}{5} =
        \frac{\sqrt 5 - 1}{4}$。
\end{enumerate}
\end{solution}

\begin{solution}{exo:b1:complex:9}
展开两个平方，与
\cref{prop:b1:complex:rules}~(4) 中的证明相同：
\[
\abs{z + w}^2 = \abs z^2 + \abs w^2 + 2\Re(z\conj w),
\qquad
\abs{z - w}^2 = \abs z^2 + \abs w^2 - 2\Re(z\conj w),
\]
两式相加。几何上，$\abs{z+w}$ 与 $\abs{z-w}$ 分别是
顶点为 $0, z, z+w, w$ 的平行四边形的两条对角线的长度，
而 $\abs z$ 与 $\abs w$ 为边长：两条对角线的
平方和等于四条边的
平方和。
\end{solution}

\begin{solution}{exo:b1:complex:10}
记 $j = \eu^{2\iu\pi/3}$，则 $j^2 + j + 1 = 0$，且 $\eu^{\iu\pi/3} =
-j^2$，$\eu^{-\iu\pi/3} = -j$。三角形为正向正三角形，当且仅当
$c - a = -j^2 (b - a)$；反向时当且仅当 $c - a = -j(b - a)$
（见\cref{ex:b1:complex:equilateral}）。

考虑 $P = a + jb + j^2 c$ 与 $Q = a + j^2 b + jc$。利用 $1 +
j^2 = -j$ 及 $j^3 = 1$：
\[
c - a + j^2(b - a) = c + j^2 b - (1 + j^2)\,a = c + j^2 b + ja
= j\,(a + jb + j^2 c) = jP,
\]
于是正向情形化为 $jP = 0 \iff P = 0$；同样计算
以 $j$ 代替 $j^2$ 可得 $c - a + j(b - a) = j^2 Q$，故
反向情形化为 $Q = 0$。因此：正三角形（不论定向）
$\iff PQ = 0$。展开，并利用 $j + j^2 = -1$：
\[
PQ = a^2 + b^2 + c^2 + (j + j^2)(ab + bc + ca)
= a^2 + b^2 + c^2 - (ab + bc + ca).
\]
故三角形为正三角形当且仅当 $a^2 + b^2 + c^2 = ab +
bc + ca$。
\end{solution}

\begin{solution}{exo:b1:complex:11}
由于 $\omega^k$、$k = 0, \dots, n-1$ 恰为 $n$ 个
$X^n - 1$ 的根（见\cref{thm:b1:complex:roots}），且
该多项式首一：
\[
X^n - 1 = \prod_{k=0}^{n-1} (X - \omega^k)
= (X - 1) \prod_{k=1}^{n-1} (X - \omega^k).
\]
两边除以 $X - 1$：$\;1 + X + \dots + X^{n-1} = \prod_{k=1}^{n-1}
(X - \omega^k)$。令 $X = 1$，得 $P = n$。

又 $1 - \omega^k = -\eu^{\iu k\pi/n}\bigl(\eu^{\iu k\pi/n} -
\eu^{-\iu k\pi/n}\bigr) = -2\iu\,\eu^{\iu k\pi/n} \sin\frac{k\pi}{n}$，
于是对 $P = n$ 取模（各 $\sin\frac{k\pi}n > 0$，当 $1 \leq k
\leq n-1$）：
\[
n = \abs P = \prod_{k=1}^{n-1} 2\sin\frac{k\pi}{n}
= 2^{n-1} \prod_{k=1}^{n-1} \sin\frac{k\pi}{n},
\qquad\text{因此}\qquad
\prod_{k=1}^{n-1} \sin\frac{k\pi}{n} = \frac{n}{2^{n-1}} .
\]
\end{solution}

\begin{solution}{exo:b1:complex:12}
$z = 1$ 不是解（$2^n \neq 0$），故方程等价于
$\bigl(\frac{z+1}{z-1}\bigr)^n = 1$，即\
$\frac{z+1}{z-1} = \omega^k$，其中 $\omega = \eu^{2\iu\pi/n}$，且
$k \in \{0, \dots, n-1\}$。$k = 0$ 应舍去（$z + 1 =
z - 1$ 不可能）。对 $1 \leq k \leq n - 1$，解
$z + 1 = \omega^k(z - 1)$ 得 $z(1 - \omega^k) = -1 - \omega^k$，
于是令 $\varphi = \frac{2k\pi}{n}$，并利用
半角分解 $1 + \eu^{\iu\varphi} = 2\cos\frac\varphi2\,
\eu^{\iu\varphi/2}$，$\eu^{\iu\varphi} - 1 = 2\iu\sin\frac\varphi2
\,\eu^{\iu\varphi/2}$：
\[
z_k = \frac{1 + \omega^k}{\omega^k - 1}
= \frac{2\cos\frac{k\pi}{n}}{2\iu\,\sin\frac{k\pi}{n}}
= -\iu\,\frac{\cos\frac{k\pi}{n}}{\sin\frac{k\pi}{n}} ,
\]
所断言的纯虚性由此成立。\omterm{def:b1:logic:map}{映射} $t \mapsto \cos t/\sin t$
在 $\intoo0\pi$ 上为\omterm{def:b1:logic:inj}{单射}（它严格递减），故
$n - 1$ 个值 $z_k$ 两两不同：方程展开后
次数为 $n - 1$（$z^n$ 项相消），恰有这
$n - 1$ 个解。
\end{solution}

\begin{solution}{pb:b1:complex:1}
\textbf{1.} $j = \cos\frac{2\pi}3 + \iu\sin\frac{2\pi}3 = -\frac12
+ \iu\frac{\sqrt3}2$; $j^2 = \eu^{4\iu\pi/3} = -\frac12 -
\iu\frac{\sqrt3}2 = \conj j$；$j^3 = 1$；$1 + j + j^2 = 0$（
三次\omterm{def:b1:complex:unity}{单位根}之和，\cref{prop:b1:complex:sumroots}）；$j^{-1}
= j^2$（因 $j \cdot j^2 = 1$）。单位圆上，$1$、$j$、
$j^2$ 是一个正向等边三角形的顶点。

\textbf{2.} 中心为 $a$、转角为 $\theta$ 的旋转保持 $a$ 不动，
并将从 $a$ 出发的每个向量旋转 $\theta$：$r(z) - a =
\eu^{\iu\theta}(z - a)$，即 $r(z) = a + \eu^{\iu\theta}(z - a)$。
复合 $r(z) = a + \eu^{\iu\theta}(z-a)$ 与 $r'(z) = a' +
\eu^{\iu\theta'}(z-a')$：
\[
r' \circ r\,(z) = \eu^{\iu(\theta + \theta')} z +
\text{常数},
\]
形如 $Az + B$ 的\omterm{def:b1:logic:map}{映射}，其中 $A = \eu^{\iu(\theta+\theta')}$ 的
模为 $1$。由\cref{prop:b1:complex:similitude}可知，它是
转角为 $\arg A = \theta + \theta'$ 的旋转（若 $A \neq 1$），而若 $A = 1$ 则为平移，
即 $\theta + \theta' \in 2\pi\Z$ 时。

\textbf{3.} $(b, c, p)$ 为等边三角形，当且仅当 $\abs{p - b} = \abs{c -
b} = \abs{p - c}$。记 $q = \frac{p - b}{c - b}$，则第一个等式
给出 $\abs q = 1$，第二个给出 $\abs{q - 1} = 1$；
合起来，$q = \eu^{\pm\iu\pi/3}$（两圆 $\abs q = 1$ 与
$\abs{q - 1} = 1$ 的两个交点为 $\frac12 \pm
\iu\frac{\sqrt3}2$）。于是 $p = b + \eu^{\pm\iu\pi/3}(c - b)$。对
$a = 0$，$b = 1$，$c = \iu$（正向三角形）：选取
$\eu^{-\iu\pi/3}$ 时得
\[
p = 1 + \Bigl(\tfrac12 - \iu\tfrac{\sqrt3}2\Bigr)(\iu - 1)
= \tfrac{1 + \sqrt3}2\,(1 + \iu) \approx 1.37 + 1.37\,\iu ,
\]
该点位于直线 $BC$（$x + y = 1$）的另一侧，与 $a
= 0$ 异侧：向外。取 $\eu^{+\iu\pi/3}$ 时得 $p \approx -0.37
- 0.37\,\iu$，与 $a$ 同侧：向内。

\textbf{4.} 将 $P = a + jb + j^2c$ 乘以 $j^2$：$j^2 P = j^2 a +
j^3 b + j^4 c = b + jc + j^2 a$（用 $j^3 = 1$，$j^4 = j$）；并
乘以 $j$ 得 $jP = ja + j^2 b + c$。这便是两个恒等式。若 $P
= 0$，则 $j^2 P = jP = 0$：判别准则对 $(b, c, a)$ 及
$(c, a, b)$ 同样成立——轮换不变性（这是理所应当的：
等边三角形并不在意先列出哪个顶点）。

\textbf{5.} 若 $a = b$：$0 = a(1 + j) + j^2 c = -j^2 a + j^2 c$
（用 $1 + j = -j^2$），故 $c = a$。若 $b = c$：$0 = a + b(j + j^2)
= a - b$，故 $a = b$。若 $a = c$：$0 = a(1 + j^2) + jb = -ja + jb$，
故 $a = b$。三种情形下三点均重合。

\textbf{6.} $(a'z + b') \circ (az + b) = a'a\,z + (a'b + b')$，其中
$a'a \neq 0$，形式相同。$z \mapsto az + b$ 的逆为 $z
\mapsto \frac1a z - \frac ba$，同样是这种形式。以恒同\omterm{def:b1:logic:map}{映射}为
单位元，所有正向相似变换在复合下构成
一个群。

\textbf{7.} $f(z) = az + b$ 满足 $f(z_1) = w_1$，$f(z_2) =
w_2$ 当且仅当 $a z_1 + b = w_1$ 且 $a z_2 + b = w_2$；相减得
$a(z_2 - z_1) = w_2 - w_1$，故
\[
a = \frac{w_2 - w_1}{z_2 - z_1} \;(\neq 0), \qquad
b = w_1 - a z_1
\]
此时系数被唯一确定，反之此取法确实可行：存在性与
唯一性。

\textbf{8.} $\abs{f(z) - f(w)} = \abs{a(z - w)} = \abs a\,\abs{z -
w}$：所有距离均乘以 $\abs a$。关于形状：
\[
\frac{f(c') - f(a')}{f(b') - f(a')} = \frac{a(c' -
a')}{a(b' - a')} = \frac{c' - a'}{b' - a'} .
\]
若两个三角形 $(a', b', c')$ 与 $(a'', b'', c'')$ 的
形状相同，设 $f$ 为满足 $f(a') = a''$
且 $f(b') = b''$（第 7 题）；则 $(a'', b'',
f(c'))$ 的形状与 $(a', b', c')$ 相同，从而与 $(a'', b'',
c'')$，而形状由前两个顶点决定第三个顶点，
即 $f(c') = c''$。反之，形状相同可由所示的
不变性得出。

\textbf{9.} $b = f(0) = 1$；$a + b = f(1) = \iu$ 给出 $a = \iu -
1$。相似比 $\abs a = \sqrt2$，转角 $\arg(\iu - 1) =
\frac{3\pi}4$。不动点：
\[
\zeta = \frac{b}{1 - a} = \frac1{2 - \iu} = \frac{2 + \iu}5 .
\]
因此 $f$ 是中心为 $\frac{2+\iu}5$ 的正向相似，相似比
为 $\sqrt2$，转角为 $\frac{3\pi}4$。

\textbf{10.} 设 $\alpha + \beta = 1$，$f(z) = az + b$：
\[
f(\alpha a' + \beta b') = a\alpha a' + a\beta b' + b
= \alpha(a a' + b) + \beta(a b' + b)
= \alpha f(a') + \beta f(b') ,
\]
关键在于 $b = (\alpha + \beta)b$。中点（$\alpha = \beta =
\frac12$）、重心（迭代可得）以及拿破仑中心 $\alpha b
+ \beta c$（下文）因此具有\emph{等变性}：对三角形作变换时，
它们也随之作相应的变换。因此，证明一个相似不变的
命题对某一个摆放恰当的三角形成立，即可推出它对所有三角形成立
——这正是第 21 题所用的依据。

\textbf{11.} $[b, c]$ 上向外的顶点为 $p_a = b +
\eu^{-\iu\pi/3}(c - b)$（第 3 题），故中心为
\[
n_a = \frac{b + c + p_a}3
= \frac{2b + c + \frac{1 - \iu\sqrt3}2\,(c - b)}3
= \frac{(3 + \iu\sqrt3)\,b + (3 - \iu\sqrt3)\,c}6
= \alpha b + \beta c ,
\]
其中 $\alpha = \frac{3 + \iu\sqrt3}6$，$\beta = \frac{3 -
\iu\sqrt3}6 = \conj\alpha$。对于边
$[c, a]$ 与 $[a, b]$ 作同样的计算，得 $n_b = \alpha c + \beta a$ 与 $n_c =
\alpha a + \beta b$（角色依次轮换）。

\textbf{12.} $\alpha + \beta = \frac{3 + \iu\sqrt3 + 3 -
\iu\sqrt3}6 = 1$。于是
\[
\frac{n_a + n_b + n_c}3
= \frac{(\alpha + \beta)(a + b + c)}3 = \frac{a + b + c}3 :
\]
两个三角形有相同的重心。

\textbf{13.} 按 $\alpha$ 与 $\beta$ 分组，并应用第 4 题的
恒等式于 $P = a + jb + j^2c$：
\[
n_a + j n_b + j^2 n_c
= \alpha\,(b + jc + j^2 a) + \beta\,(c + ja + j^2 b)
= \alpha\,j^2 P + \beta\,j P
= (\alpha j^2 + \beta j)\,P .
\]

\textbf{14.} 取 $j = \frac{-1 + \iu\sqrt3}2$：
\[
\alpha j = \frac{(3 + \iu\sqrt3)(-1 + \iu\sqrt3)}{12}
= \frac{-3 + 3\iu\sqrt3 - \iu\sqrt3 - 3}{12}
= \frac{-3 + \iu\sqrt3}6 = -\beta ,
\]
故 $\alpha j + \beta = 0$，从而 $\alpha j^2 + \beta j = j(\alpha j
+ \beta) = 0$，由第 13 题得 $n_a + jn_b + j^2n_c = 0$，
对每个三角形 $(a, b, c)$ 成立。由判别准则（第 4--5 题），
这些中心构成正向等边三角形，或为单个点：
拿破仑定理。

\textbf{15.} 向内的顶点为 $b + \eu^{+\iu\pi/3}(c - b)$，
并沿用第 11 题的计算，取 $\eu^{+\iu\pi/3} = \frac{1 +
\iu\sqrt3}2$，则 $\alpha$ 与 $\beta$ 互换：$m_a = \beta b + \alpha
c$，$m_b = \beta c + \alpha a$，$m_c = \beta a + \alpha b$。利用
类似的恒等式 $b + j^2c + ja = j\,Q$ 与 $c + j^2a + jb
= j^2 Q$，其中 $Q = a + j^2b + jc$：
\[
m_a + j^2 m_b + j m_c
= \beta(b + j^2 c + ja) + \alpha(c + j^2 a + jb)
= (\beta j + \alpha j^2)\, Q = j(\beta + \alpha j)\,Q = 0 ,
\]
由于 $\alpha j = -\beta$（第 14 题），因此内侧的中心
满足\emph{反向}等边判别准则，即
为方向相反的等边三角形（或为一个点）。

\textbf{16.} 对三元组 $(n_a, n_b, n_c)$ 应用第 5 题，
（它满足正向判别准则），两个中心重合当且仅当
三个中心全部重合；而三个中心全部重合当且仅当\emph{两个}准则同时成立，
即另须 $n_a + j^2 n_b + j n_c = 0$。
仿第 13 题计算，利用恒等式 $b + j^2c + ja = jQ$，$c + j^2a +
jb = j^2Q$：
\[
n_a + j^2 n_b + j n_c
= \alpha\,jQ + \beta\,j^2 Q = j(\alpha + \beta j)\,Q .
\]
直接计算：$\beta j = \frac{(3 -
\iu\sqrt3)(-1 + \iu\sqrt3)}{12} = \frac{-3 + 3\iu\sqrt3 +
\iu\sqrt3 + 3}{12} = \frac{\iu\sqrt3}3$，故 $\alpha + \beta j =
\frac{3 + \iu\sqrt3 + 2\iu\sqrt3}6 = \frac{1 + \iu\sqrt3}2 \neq
0$。由此，外侧拿破仑三角形退化当且仅当 $Q = a +
j^2b + jc = 0$，即 $(a, b, c)$ 为
反向等边三角形——此时“向外”的作法全都指向
三角形外接圆所围区域内，并共享同一个中心。

\textbf{17.} 取 $a = 0$，$b = 1$，$c = \iu$：
\[
n_a = \alpha + \beta\iu = \frac{(3 + \sqrt3)(1 + \iu)}6,
\qquad
n_b = \alpha\iu = \frac{-\sqrt3 + 3\iu}6,
\qquad
n_c = \beta = \frac{3 - \iu\sqrt3}6 .
\]
边长的平方：$n_a - n_b = \frac{(3 + 2\sqrt3) +
\iu\sqrt3}6$，得 $\abs{n_a - n_b}^2 = \frac{(3 + 2\sqrt3)^2 +
3}{36} = \frac{24 + 12\sqrt3}{36} = \frac{2 + \sqrt3}3$; $n_b -
n_c = \frac{(3 + \sqrt3)(-1 + \iu)}6$，得 $\abs{n_b - n_c}^2 =
\frac{2(3 + \sqrt3)^2}{36} = \frac{24 + 12\sqrt3}{36}$; $n_c - n_a
= \frac{-\sqrt3 - \iu(3 + 2\sqrt3)}6$ 同样得到相同的值。
三边长的平方均为 $\frac{2 + \sqrt3}3$，故为等边三角形，
正如所言。

\textbf{18.} 展开：
\[
(a - b)(c - d) + (a - d)(b - c)
= (ac - ad - bc + bd) + (ab - ac - bd + cd)
= ab - ad - bc + cd ,
\]
而 $(a - c)(b - d) = ab - ad - bc + cd$，二者相等。

\textbf{19.} 对第 18 题两边取模，并应用
三角不等式（\cref{prop:b1:complex:rules}~(4)）：
\[
AC \cdot BD = \abs{(a-b)(c-d) + (a-d)(b-c)}
\leq \abs{a-b}\,\abs{c-d} + \abs{a-d}\,\abs{b-c}
= AB \cdot CD + AD \cdot BC ,
\]
等号成立当且仅当两加项位于从 $0$ 出发的
同一条射线上（即三角不等式的取等情形）。

\textbf{20.} 半角因式分解
（\cref{met:b1:complex:trig}~(4)）：
$\eu^{\iu\theta} - \eu^{\iu\varphi} = 2\iu\,\sin\frac{\theta -
\varphi}2\;\eu^{\iu(\theta + \varphi)/2}$，取模即可消去
单位模因子：$\abs{\eu^{\iu\theta} - \eu^{\iu\varphi}}
= 2\,\abs{\sin\frac{\theta - \varphi}2}$.

\textbf{21.} 取 $p = \eu^{\iu\theta}$，$\theta \in
\intoo{2\pi/3}{4\pi/3}$，由第 20 题得
\[
PA = 2\,\abs{\sin\tfrac\theta2},
\quad
PB = 2\,\abs{\sin\bigl(\tfrac\theta2 - \tfrac\pi3\bigr)},
\quad
PC = 2\,\abs{\sin\bigl(\tfrac\theta2 - \tfrac{2\pi}3\bigr)} .
\]
在该弧上，$\frac\theta2 \in \intoo{\pi/3}{2\pi/3}$：于是
$\sin\frac\theta2 > 0$; $\frac\theta2 - \frac\pi3 \in
\intoo0{\pi/3}$，故第二个正弦为正；$\frac\theta2 -
\frac{2\pi}3 \in \intoo{-\pi/3}0$，故第三个为负，
从而 $PC = 2\sin\bigl(\frac{2\pi}3 - \frac\theta2\bigr)$。和差
化积：
\[
\sin\Bigl(\frac\theta2 - \frac\pi3\Bigr) +
\sin\Bigl(\frac{2\pi}3 - \frac\theta2\Bigr)
= 2\,\sin\frac\pi6\,\cos\Bigl(\frac\theta2 - \frac\pi2\Bigr)
= \sin\frac\theta2 ,
\]
故 $PB + PC = 2\sin\frac\theta2 = PA$：范斯霍滕定理。

\textbf{22.} 当 $p = -1$ 时：$PA = 2$，$PB = PC = 1$，且
正三角形的各边长均为 $\sqrt3$，故托勒密等式两边都为
$2\sqrt3$。代数上，利用 $-1 - j^2 = j$ 及
$1 + j = -j^2$：
\[
(a - b)(p - c) = (1 - j)\,(-1 - j^2) = (1 - j)j = j - j^2
= \iu\sqrt3 ,
\]
\[
(a - c)(b - p) = (1 - j^2)(j + 1) = 1 + j - j^2 - j^3
= j - j^2 = \iu\sqrt3 .
\]
两项相等，故其比为
$1 \in \R_{>0}$：即第 19 题中的等号情形，恰与 $PA
\cdot BC = PB \cdot AC + PC \cdot AB$ 完全吻合。

\textbf{23.} 取 $(a, b, c) = (1, j, j^2)$，则 $n_a = \alpha j +
\beta j^2 = -\beta + \beta j^2$（第 14 题）$= \beta(j^2 - 1)$；
数值上 $\beta(j^2 - 1) = \frac{(3 - \iu\sqrt3)}6 \cdot
\bigl(-\frac32 - \iu\frac{\sqrt3}2\bigr) = -1$。同理 $n_b =
\alpha j^2 + \beta = -j$，且 $n_c = \alpha + \beta j = -j^2$。
$(1, j, j^2)$ 的外拿破仑三角形为 $(-1, -j, -j^2)$：
即原三角形关于其重心 $0$ 作中心对称——大小相同，
仅相差半圈旋转。正三角形是拿破仑作图中的
一个不变形状，而非收缩的极限。

\textbf{24.} (i) 乘以模为 1 的数，作为旋转，
构造了顶点与中心（第 2--3、11 题），并将
圆上的距离化为正弦（第 20 题）。(ii) 群
结构与形状不变量证明了将
外接圆归一化为单位圆、将三角形归一化为 $(1, j, j^2)$
（第 21 题）中，依据第 10 题的等变性。(iii) 半角分解
同时支撑了第 20 题与
完成范斯霍滕定理的和差化积一步。

\textbf{25.} 该翻译规则把几何命题化为多项式
恒等式，其中每个假设都是一个等式：
其收获是机械化——拿破仑定理被化为
$\alpha j + \beta = 0$，
在 $\Q(\iu\sqrt3)$ 中只需一行算术；其代价是几何直观性——该计算
只能验证结论，却不能说明\emph{为何}三个中心
恰好构成正三角形。一个较为公允的回答是，
该恒等式解释了定理的\emph{必然性}（它对
$a, b, c$ 恒成立，无需构形上的巧思，
），而图形解释其\emph{内容}。同一规则
以矩阵形式再次出现，即旋转化为
由\cref{ch:b1:matrices}与\cref{ch:b1:euclid}给出的正交 $2 \times 2$ 矩阵时，
其中“乘以
$\eu^{\iu\theta}$”是线性等距变换的原型。
\end{solution}
